% chapters/05_ergebnisse_und_diskussion.tex
%
% Erstentwurf vom 27.10.2026 (KI-gestützt erstellt, siehe Erklärung).
% Struktur, Tabellenskelette, Konventionen und Deutungsrahmen stehen;
% alle Messwerte und alle Aussagen, deren Richtung von den Messwerten
% abhängt, sind mit \PH{...} rot markiert und aus den Laufartefakten
% (Anhang~A.1: evaluation_metrics.yaml, inference_metrics.yaml,
% test_predictions.json) zu füllen. Die "% TODO"-Kommentare je Abschnitt
% nennen die dafür nötigen Auswertungen.
%
% Konventionen dieses Kapitels:
%  - Alle F1-, Precision- und Recall-Werte in Punkten (0--100).
%  - Trainierte Konfigurationen: Mittelwert +- Stichprobenstandardabweichung
%    über die Seeds 42, 123, 456; Zero-Shot: Einzellauf (Seed 42).
%  - Abschnitte 5.1 bis 5.6 berichten, 5.7 interpretiert.
%  - Bereits bekannte Werte (Support, Parameterzahlen, GPU-Typen,
%    gewählte Epochen, 8,7 ms Encoder-Latenz) sind eingetragen.
%
% Hinweis für Kapitel 4: Dort steht in 4.4.3 "strikter F1-Wert unter 0,01"
% (Skala 0--1). Für Konsistenz mit diesem Kapitel auf "unter 1" ändern.
\providecommand{\PH}[1]{\textcolor{red}{\textbf{[#1]}}}

\chapter{Ergebnisse und Diskussion}
\label{chap:ergebnisse}

Dieses Kapitel berichtet die Ergebnisse der in Kapitel~\ref{chap:methodik}
beschriebenen Experimente und diskutiert sie anschließend.
Abschnitt~\ref{sec:erg-leistung} stellt die Erkennungsleistung aller
Konfigurationen insgesamt und je Entitätstyp dar,
Abschnitt~\ref{sec:erg-formtreue} die Formtreue der LLM-Ausgaben; beide
zusammen beantworten FF1.
Abschnitt~\ref{sec:erg-groesse} betrachtet den Einfluss der Modellgröße,
Abschnitt~\ref{sec:erg-effizienz} den Ressourcenbedarf (FF2).
Abschnitt~\ref{sec:erg-fehlertypen} vergleicht die Fehlerprofile der
Modellklassen, Abschnitt~\ref{sec:erg-crossner} berichtet den
CrossNER-Ergänzungsversuch.
Abschnitt~\ref{sec:erg-diskussion} ordnet die Befunde in die verwandten
Arbeiten ein, leitet daraus die Empfehlungen zu FF3 ab und benennt,
welche Unterschiede angesichts der Streuung belastbar sind.
Die Abschnitte~\ref{sec:erg-leistung} bis~\ref{sec:erg-crossner}
beschränken sich auf die Darstellung der Messwerte; Interpretationen
folgen in Abschnitt~\ref{sec:erg-diskussion}.

Für alle Kennzahlen gelten die Festlegungen aus
Abschnitt~\ref{sec:evaluationskonzept}.
Precision, Recall und F$_1$ sind strikte Entity-Level-Werte nach dem
IOB2-Schema, berechnet mit \texttt{seqeval} auf dem vollständigen
englischen MultiNERD-Testsplit mit 16\,454 Sätzen und 25\,645
Referenzentitäten (Tabelle~\ref{tab:datensplits}); sie werden in Punkten
auf einer Skala von 0 bis 100 angegeben.
Für die fünf trainierten Konfigurationen werden Mittelwert und
Stichprobenstandardabweichung über die Seeds 42, 123 und 456 berichtet,
in den Tabellen als Mittelwert~$\pm$~Standardabweichung.
Die drei Zero-Shot-Konfigurationen beruhen auf je einem Lauf mit
Sampling (Abschnitt~\ref{sec:llm-pipeline}); ihre Werte sind daher
Einzelwerte ohne Streuungsangabe.
Ein Unterschied zwischen zwei trainierten Konfigurationen wird im
Folgenden nur dann als Unterschied bezeichnet, wenn er größer ist als
die Summe der beiden Standardabweichungen; kleinere Differenzen werden
als nicht unterscheidbar bezeichnet.
Dieses Kriterium ist in Abschnitt~\ref{sec:performanzmetriken}
festgelegt; es ersetzt keinen statistischen Test, macht die Lesart aber
nachvollziehbar.

\section{Erkennungsleistung auf MultiNERD}
\label{sec:erg-leistung}

Dieser Abschnitt berichtet die Erkennungsleistung zunächst über alle
Typen hinweg (Abschnitt~\ref{sec:erg-gesamt}) und anschließend je
Entitätstyp (Abschnitt~\ref{sec:erg-typen}).

\subsection{Gesamtergebnisse}
\label{sec:erg-gesamt}

% TODO: Werte aus evaluation_metrics.yaml der sechs DeBERTa-Läufe:
%  - Precision, Recall, Micro-F1 je Seed -> Mittelwert und Stichproben-SD
%  - Makro-F1 aus den typbezogenen F1-Werten (ungewichtetes Mittel über 15 Typen)
%  - Seed-Einzelwerte in tab:app-seeds (Anhang A.2)

Tabelle~\ref{tab:haupt} fasst die Erkennungsleistung aller acht
Konfigurationen zusammen; die Seed-Einzelwerte stehen in
Anhang~\ref{app:ergebnistabellen}.
Die beiden Encoder stehen in den ersten beiden Zeilen.

\begin{table}[htbp]
  \centering
  \small
  \caption{Erkennungsleistung aller Konfigurationen auf dem englischen
  MultiNERD-Testsplit (strikte Entity-Level-Werte, IOB2, in Punkten).
  Trainierte Konfigurationen: Mittelwert~$\pm$~Stichprobenstandardabweichung
  über drei Seeds; Zero-Shot: Einzellauf mit Sampling.
  Makro-F$_1$: ungewichtetes Mittel der typbezogenen F$_1$-Werte über
  15 Typen.
  Seed-Einzelwerte in Anhang~\ref{app:ergebnistabellen}.}
  \label{tab:haupt}
  \setlength{\tabcolsep}{4pt}
  \begin{tabular}{llcccc}
    \toprule
    \textbf{Modell} & \textbf{Regime} & \textbf{Precision} &
    \textbf{Recall} & \textbf{Micro-F$_1$} & \textbf{Makro-F$_1$} \\
    \midrule
    DeBERTa-v3-base  & Fine-Tuning & \PH{xx,x}\,$\pm$\,\PH{x,x} & \PH{xx,x}\,$\pm$\,\PH{x,x} & \PH{xx,x}\,$\pm$\,\PH{x,x} & \PH{xx,x} \\
    DeBERTa-v3-large & Fine-Tuning & \PH{xx,x}\,$\pm$\,\PH{x,x} & \PH{xx,x}\,$\pm$\,\PH{x,x} & \PH{xx,x}\,$\pm$\,\PH{x,x} & \PH{xx,x} \\
    \midrule
    Qwen3.5-0.8B & Zero-Shot & \PH{xx,x} & \PH{xx,x} & \PH{xx,x} & \PH{xx,x} \\
    Qwen3.5-4B   & Zero-Shot & \PH{xx,x} & \PH{xx,x} & \PH{xx,x} & \PH{xx,x} \\
    Qwen3.5-27B  & Zero-Shot & \PH{xx,x} & \PH{xx,x} & \PH{xx,x} & \PH{xx,x} \\
    \midrule
    Qwen3.5-0.8B & QLoRA & \PH{xx,x}\,$\pm$\,\PH{x,x} & \PH{xx,x}\,$\pm$\,\PH{x,x} & \PH{xx,x}\,$\pm$\,\PH{x,x} & \PH{xx,x} \\
    Qwen3.5-4B   & QLoRA & \PH{xx,x}\,$\pm$\,\PH{x,x} & \PH{xx,x}\,$\pm$\,\PH{x,x} & \PH{xx,x}\,$\pm$\,\PH{x,x} & \PH{xx,x} \\
    Qwen3.5-27B  & QLoRA & \PH{xx,x}\,$\pm$\,\PH{x,x} & \PH{xx,x}\,$\pm$\,\PH{x,x} & \PH{xx,x}\,$\pm$\,\PH{x,x} & \PH{xx,x} \\
    \bottomrule
  \end{tabular}
\end{table}

DeBERTa-v3-base erreicht einen Micro-F$_1$ von \PH{xx,x}, DeBERTa-v3-large
einen von \PH{xx,x}.
Die Differenz von \PH{x,x} Punkten ist \PH{größer / kleiner} als die Summe
der beiden Standardabweichungen (\PH{x,x} und \PH{x,x}); die größere
Variante ist damit \PH{unterscheidbar besser / nicht unterscheidbar}.
Precision und Recall liegen bei beiden Modellen \PH{nahe beieinander /
auseinander}: \PH{Precision xx,x gegenüber Recall xx,x (base) bzw.
xx,x gegenüber xx,x (large)}.
Der Makro-F$_1$ liegt mit \PH{xx,x} bzw. \PH{xx,x} um \PH{x} Punkte unter
dem Micro-F$_1$, was auf \PH{deutlich schwächere Werte bei den seltenen
oder verrauschten Typen} zurückgeht (Tabelle~\ref{tab:typen} in
Abschnitt~\ref{sec:erg-typen}).

Die Streuung über die Seeds ist mit Standardabweichungen von \PH{x,x}
(base) und \PH{x,x} (large) \PH{gering}.
Die Checkpoint-Auswahl fiel bei DeBERTa-v3-base auf die vierte
(Seed~42) bzw. fünfte Epoche (Seeds 123 und 456), bei DeBERTa-v3-large
auf die dritte bzw. vierte Epoche (Abschnitt~\ref{sec:encoder-pipeline});
die Validierungs-F$_1$-Werte der ausgewählten Checkpoints liegen bei
\PH{xx,x} bzw. \PH{xx,x} und damit \PH{über / nahe an} den Testwerten.
% TODO: Validierungs-F1 der gewählten Checkpoints aus results.yaml;
% Abweichung Validierung/Test kommentieren (Verteilungsverschiebung LOC
% 29,1 % -> 46,9 %, Abschnitt 4.3.3).


% TODO: Werte aus den drei Zero-Shot-Läufen (inference_metrics.yaml bzw.
% aus test_predictions.json nachgerechnet) und den neun QLoRA-Läufen
% (evaluation_metrics.yaml). Parse-Fehlerraten gehören in tab:formtreue-erg (Abschnitt 5.2).

Die Zeilen drei bis acht in Tabelle~\ref{tab:haupt} zeigen die
LLM-Konfigurationen.
Im Zero-Shot-Regime erreichen die drei Qwen3.5-Modelle Micro-F$_1$-Werte
von \PH{xx,x} (0.8B), \PH{xx,x} (4B) und \PH{xx,x} (27B).
\PH{Beschreibung des Profils: Precision gegenüber Recall je Größe, etwa
ob die Modelle viele Entitäten übersehen (niedriger Recall) oder
Nicht-Entitäten markieren (niedrige Precision).}
Die Werte sind Einzelwerte aus gesampelten Ausgaben und beruhen auf einem
Prompt und Decoding-Einstellungen, die nach Beobachtungen auf dem
Testsplit gewählt wurden (Abschnitt~\ref{sec:datenaufteilung}); sie sind
daher als obere Abschätzung dessen zu lesen, was die unveränderten
Checkpoints mit diesem Format leisten.

Nach der Anpassung mit QLoRA steigt der Micro-F$_1$ auf \PH{xx,x}
(0.8B), \PH{xx,x} (4B) und \PH{xx,x} (27B).
Der Zugewinn gegenüber dem Zero-Shot-Regime beträgt damit \PH{xx}, \PH{xx}
bzw. \PH{xx} Punkte und ist \PH{bei allen drei Größen deutlich größer als
die Streuung über die Seeds}.
Die Standardabweichungen über die Seeds liegen bei \PH{x,x}, \PH{x,x} und
\PH{x,x}.
Der Makro-F$_1$ der QLoRA-Modelle beträgt \PH{xx,x}, \PH{xx,x} und
\PH{xx,x}.
Ausgewählt wurde in acht der neun Läufe der Adapter der dritten Epoche,
beim 4B-Modell mit Seed~42 der der zweiten
(Abschnitt~\ref{sec:llm-pipeline}); die generative Validierung auf 200
Sätzen ergab für die gewählten Adapter F$_1$-Werte von \PH{xx,x},
\PH{xx,x} und \PH{xx,x}.
% TODO: Validierungs-F1 (200 Sätze, greedy) aus results.yaml; prüfen, ob
% der Verlauf über die Epochen (Epoche 1/2/3) bereits sättigt.


Der beste Encoder, \PH{DeBERTa-v3-large}, erreicht einen Micro-F$_1$ von
\PH{xx,x}; das beste LLM, \PH{Qwen3.5-27B im QLoRA-Regime}, einen von
\PH{xx,x}.
Der Abstand beträgt \PH{x,x} Punkte zugunsten \PH{des Encoders / des LLM}
und ist \PH{größer / kleiner} als die Summe der Standardabweichungen.
\PH{Einordnung der kleineren Konfigurationen: Liegt Qwen3.5-4B (QLoRA)
über oder unter DeBERTa-v3-base? Wo steht Qwen3.5-0.8B (QLoRA) relativ zu
den Encodern?}
Alle Zero-Shot-Konfigurationen liegen \PH{xx bis xx} Punkte unter dem
schwächsten trainierten Modell.

Die beiden Modellklassen unterscheiden sich im Profil aus Precision und
Recall: \PH{Encoder: Precision xx,x / Recall xx,x; bestes LLM: Precision
xx,x / Recall xx,x. Beschreiben, ob der Vorsprung einer Klasse vor allem
aus der Precision oder aus dem Recall stammt.}

\subsection{Ergebnisse nach Entitätstyp}
\label{sec:erg-typen}


% TODO: typbezogene P/R/F1 je Konfiguration (evaluation_metrics.yaml für
% die trainierten Läufe; Zero-Shot aus test_predictions.json nachrechnen).
% Mittelwerte über Seeds; vollständige Tabelle mit P/R/F1 in Anhang A.2.

\begin{table}[htbp]
  \centering
  \footnotesize
  \caption{Typbezogene F$_1$-Werte (strikt, in Punkten) ausgewählter
  Konfigurationen auf dem Testsplit; trainierte Konfigurationen als
  Mittelwert über drei Seeds, Zero-Shot als Einzellauf.
  Support: Zahl der Referenzentitäten im Testsplit
  (Tabelle~\ref{tab:datensplits}).
  Die Typen sind nach der Zuverlässigkeit ihrer Annotation gruppiert
  (Abschnitt~\ref{sec:datenauswahl}).
  Precision und Recall je Typ in Anhang~\ref{app:ergebnistabellen}.}
  \label{tab:typen}
  \setlength{\tabcolsep}{3pt}
  \begin{tabular}{lrcccccc}
    \toprule
    & & \multicolumn{2}{c}{\textbf{DeBERTa-v3}} &
      \multicolumn{3}{c}{\textbf{Qwen3.5 QLoRA}} &
      \textbf{Qwen3.5} \\
    \cmidrule(lr){3-4}\cmidrule(lr){5-7}\cmidrule(lr){8-8}
    \textbf{Typ} & \textbf{Support} & \textbf{base} & \textbf{large} &
      \textbf{0.8B} & \textbf{4B} & \textbf{27B} & \textbf{27B Zero-Shot} \\
    \midrule
    \multicolumn{8}{l}{\emph{Zuverlässig annotierte Typen}} \\
    PER   & 5\,265  & \PH{xx,x} & \PH{xx,x} & \PH{xx,x} & \PH{xx,x} & \PH{xx,x} & \PH{xx,x} \\
    ORG   & 3\,309  & \PH{xx,x} & \PH{xx,x} & \PH{xx,x} & \PH{xx,x} & \PH{xx,x} & \PH{xx,x} \\
    LOC   & 12\,024 & \PH{xx,x} & \PH{xx,x} & \PH{xx,x} & \PH{xx,x} & \PH{xx,x} & \PH{xx,x} \\
    CEL   & 41      & \PH{xx,x} & \PH{xx,x} & \PH{xx,x} & \PH{xx,x} & \PH{xx,x} & \PH{xx,x} \\
    DIS   & 759     & \PH{xx,x} & \PH{xx,x} & \PH{xx,x} & \PH{xx,x} & \PH{xx,x} & \PH{xx,x} \\
    EVE   & 352     & \PH{xx,x} & \PH{xx,x} & \PH{xx,x} & \PH{xx,x} & \PH{xx,x} & \PH{xx,x} \\
    MEDIA & 458     & \PH{xx,x} & \PH{xx,x} & \PH{xx,x} & \PH{xx,x} & \PH{xx,x} & \PH{xx,x} \\
    \midrule
    \multicolumn{8}{l}{\emph{Häufig verwechselte oder unvollständig verlinkte Typen}} \\
    ANIM  & 1\,604  & \PH{xx,x} & \PH{xx,x} & \PH{xx,x} & \PH{xx,x} & \PH{xx,x} & \PH{xx,x} \\
    BIO   & 8       & \PH{xx,x} & \PH{xx,x} & \PH{xx,x} & \PH{xx,x} & \PH{xx,x} & \PH{xx,x} \\
    FOOD  & 566     & \PH{xx,x} & \PH{xx,x} & \PH{xx,x} & \PH{xx,x} & \PH{xx,x} & \PH{xx,x} \\
    PLANT & 894     & \PH{xx,x} & \PH{xx,x} & \PH{xx,x} & \PH{xx,x} & \PH{xx,x} & \PH{xx,x} \\
    \midrule
    \multicolumn{8}{l}{\emph{Übrige Typen}} \\
    INST  & 12      & \PH{xx,x} & \PH{xx,x} & \PH{xx,x} & \PH{xx,x} & \PH{xx,x} & \PH{xx,x} \\
    MYTH  & 32      & \PH{xx,x} & \PH{xx,x} & \PH{xx,x} & \PH{xx,x} & \PH{xx,x} & \PH{xx,x} \\
    TIME  & 289     & \PH{xx,x} & \PH{xx,x} & \PH{xx,x} & \PH{xx,x} & \PH{xx,x} & \PH{xx,x} \\
    VEHI  & 32      & \PH{xx,x} & \PH{xx,x} & \PH{xx,x} & \PH{xx,x} & \PH{xx,x} & \PH{xx,x} \\
    \midrule
    Micro-F$_1$ (alle) & 25\,645 & \PH{xx,x} & \PH{xx,x} & \PH{xx,x} & \PH{xx,x} & \PH{xx,x} & \PH{xx,x} \\
    Makro-F$_1$ & -- & \PH{xx,x} & \PH{xx,x} & \PH{xx,x} & \PH{xx,x} & \PH{xx,x} & \PH{xx,x} \\
    \bottomrule
  \end{tabular}
\end{table}

Tabelle~\ref{tab:typen} schlüsselt die F$_1$-Werte nach Entitätstyp auf;
Abbildung~\ref{fig:typen} stellt die drei Kernkonfigurationen gegenüber.
Bei den drei häufigen Typen PER, ORG und LOC, die 80\,\% der
Testentitäten stellen, liegen alle trainierten Konfigurationen
\PH{über xx Punkten}; die Unterschiede zwischen den Modellklassen betragen
hier \PH{höchstens x,x Punkte}.
Die Typen ANIM, BIO, FOOD und PLANT, deren Silberannotation nach
\textcite{tedeschi2022multinerd} häufig verwechselt oder unvollständig
verlinkt ist, erreichen bei allen Modellen \PH{deutlich niedrigere} Werte
zwischen \PH{xx} und \PH{xx} Punkten.
\PH{Benennen, bei welchen Typen die Modellklassen am stärksten
auseinanderliegen, etwa ob das LLM bei DIS, EVE oder MEDIA vorn liegt
und der Encoder bei LOC oder ANIM.}
Für BIO (8 Referenzentitäten), INST (12), MYTH (32) und VEHI (32) sind
die Werte mit erheblicher Unsicherheit behaftet, da wenige Fehler den
F$_1$ um zweistellige Beträge verändern; sie werden nicht einzeln
interpretiert.
Im Zero-Shot-Regime \PH{beschreiben: welche Typen das 27B-Modell ohne
Training erkennt (vermutlich PER, LOC, ORG, DIS) und welche nahezu nicht
(vermutlich MEDIA, EVE, TIME mit abweichenden Typgrenzen)}.


\begin{figure}[htbp]
  \centering
  \fbox{\parbox[c][6cm][c]{0.9\linewidth}{\centering
    \PH{Abbildung: F$_1$ je Entitätstyp (y) für DeBERTa-v3-large,
    Qwen3.5-27B (QLoRA) und Qwen3.5-27B (Zero-Shot); Typen auf der
    x-Achse nach Support sortiert, Gruppen aus Tabelle~\ref{tab:typen}
    farblich abgesetzt; ein Punkt je Konfiguration und Typ}}}
  \caption{Typbezogene F$_1$-Werte der drei Kernkonfigurationen auf dem
  Testsplit (Werte aus Tabelle~\ref{tab:typen}).}
  \label{fig:typen}
\end{figure}

\section{Formtreue der LLM-Ausgaben}
\label{sec:erg-formtreue}

% TODO: Zähler aus inference_metrics.yaml je LLM-Lauf (Kennzahlen nach
% Tab. 4.8); QLoRA als Mittel über drei Seeds, gerundet; zusätzlich die
% Zahl nicht parsebarer Antworten ohne schließende Klammer (Indikator für
% Abbruch am 256-Token-Limit) aus test_predictions.json.

Tabelle~\ref{tab:formtreue-erg} berichtet die Kennzahlen zur Formtreue
aus Tabelle~\ref{tab:formtreue} für die sechs LLM-Konfigurationen.

\begin{table}[htbp]
  \centering
  \footnotesize
  \caption{Formtreue der LLM-Ausgaben auf dem Testsplit (16\,454
  Antworten).
  Parse-Fehlerrate in Prozent der Antworten; übrige Werte absolute
  Häufigkeiten (QLoRA: Mittel über drei Seeds, gerundet).
  Jedes verworfene Element zählt nur bei der ersten verletzten Regel
  (Abschnitt~\ref{sec:robustheit-ausgaben}).
  Leere Vorhersagen: Sätze ohne übernommenen Span.}
  \label{tab:formtreue-erg}
  \setlength{\tabcolsep}{4pt}
  \begin{tabular}{>{\raggedright\arraybackslash}p{4.4cm}rrrrrr}
    \toprule
    & \multicolumn{3}{c}{\textbf{Zero-Shot}} &
      \multicolumn{3}{c}{\textbf{QLoRA}} \\
    \cmidrule(lr){2-4}\cmidrule(lr){5-7}
    \textbf{Kennzahl} & \textbf{0.8B} & \textbf{4B} & \textbf{27B} &
      \textbf{0.8B} & \textbf{4B} & \textbf{27B} \\
    \midrule
    Parse-Fehlerrate [\%]          & \PH{x,x} & \PH{x,x} & \PH{x,x} & \PH{x,x} & \PH{x,x} & \PH{x,x} \\
    \quad davon ohne schließende Klammer & \PH{x} & \PH{x} & \PH{x} & \PH{x} & \PH{x} & \PH{x} \\
    Parse direkt / Codeblock / Extraktion & \PH{x/x/x} & \PH{x/x/x} & \PH{x/x/x} & \PH{x/x/x} & \PH{x/x/x} & \PH{x/x/x} \\
    Falsche Struktur (kein Array)  & \PH{x} & \PH{x} & \PH{x} & \PH{x} & \PH{x} & \PH{x} \\
    Nicht-Objekte                  & \PH{x} & \PH{x} & \PH{x} & \PH{x} & \PH{x} & \PH{x} \\
    Fehlende Felder                & \PH{x} & \PH{x} & \PH{x} & \PH{x} & \PH{x} & \PH{x} \\
    Unzulässige Indizes            & \PH{x} & \PH{x} & \PH{x} & \PH{x} & \PH{x} & \PH{x} \\
    Unbekannte Typen               & \PH{x} & \PH{x} & \PH{x} & \PH{x} & \PH{x} & \PH{x} \\
    Text-Abweichungen              & \PH{x} & \PH{x} & \PH{x} & \PH{x} & \PH{x} & \PH{x} \\
    Überlappungen                  & \PH{x} & \PH{x} & \PH{x} & \PH{x} & \PH{x} & \PH{x} \\
    \midrule
    Leere Vorhersagen & \PH{x} & \PH{x} & \PH{x} & \PH{x} & \PH{x} & \PH{x} \\
    \bottomrule
  \end{tabular}
\end{table}

Im QLoRA-Regime ist die Parse-Fehlerrate mit \PH{x,x}, \PH{x,x} und
\PH{x,x}\,\% \PH{gering}; \PH{nahezu alle / x von x} nicht parsebaren
Antworten enden ohne schließende Klammer und gehen damit mit hoher
Wahrscheinlichkeit auf das Limit von 256 Ausgabetokens zurück.
Unbekannte Typen und Text-Abweichungen treten nach dem Training
\PH{praktisch nicht mehr / in x Fällen} auf; die häufigste verbleibende
Fehlerklasse ist \PH{Überlappungen / unzulässige Indizes}.
Im Zero-Shot-Regime zeigt sich ein anderes Bild: \PH{Beschreiben je
Größe: Parse-Fehlerrate, Anteil der Extraktion aus umgebendem Text,
halluzinierte Typen (Beispiele nennen, etwa ausgeschriebene Typnamen
oder neue Kategorien), Text-Abweichungen als Indiz für fehlerhaft
gezählte Indizes, leere Arrays.}
Da die Zähler absolute Häufigkeiten ohne Bezug zur Zahl erzeugter
Elemente sind, lassen sie sich zwischen Konfigurationen vergleichen,
aber nicht als Fehlerraten je Element lesen
(Abschnitt~\ref{sec:robustheit-ausgaben}).

\section{Einfluss der Modellgröße auf die NER-Leistung}
\label{sec:erg-groesse}

% TODO: Differenzen zwischen den Größenklassen aus tab:haupt berechnen;
% Abbildung fig:groesse (F1 gegen Parameterzahl, log-x, Fehlerbalken = SD) aus
% den Mittelwerten erzeugen (pgfplots oder extern als PDF).

Tabelle~\ref{tab:groesse} und Abbildung~\ref{fig:groesse} stellen die
Erkennungsleistung über die Größenklassen aus
Abschnitt~\ref{sec:groessenklassen} getrennt nach Familie und Regime dar.

\begin{table}[htbp]
  \centering
  \small
  \caption{Micro-F$_1$ (in Punkten) nach Größenklasse und Regime.
  S/M/L bezeichnen die relative Größe innerhalb der Familie
  (Tabelle~\ref{tab:experiment-matrix}); $\Delta$ ist die Differenz zur
  nächstkleineren Klasse.
  Trainierte Konfigurationen als Mittelwert~$\pm$~Standardabweichung über
  drei Seeds.}
  \label{tab:groesse}
  \setlength{\tabcolsep}{3pt}
  \begin{tabular}{llccccc}
    \toprule
    \textbf{Familie} & \textbf{Regime} & \textbf{S} & \textbf{M} &
    \textbf{L} & $\Delta$(M$-$S) & $\Delta$(L$-$M) \\
    \midrule
    DeBERTa-v3 & Fine-Tuning & \PH{xx,x}\,$\pm$\,\PH{x,x} & \PH{xx,x}\,$\pm$\,\PH{x,x} & -- & \PH{+x,x} & -- \\
    Qwen3.5    & Zero-Shot   & \PH{xx,x} & \PH{xx,x} & \PH{xx,x} & \PH{+x,x} & \PH{+x,x} \\
    Qwen3.5    & QLoRA       & \PH{xx,x}\,$\pm$\,\PH{x,x} & \PH{xx,x}\,$\pm$\,\PH{x,x} & \PH{xx,x}\,$\pm$\,\PH{x,x} & \PH{+x,x} & \PH{+x,x} \\
    \bottomrule
  \end{tabular}
\end{table}

\begin{figure}[htbp]
  \centering
  \fbox{\parbox[c][6cm][c]{0.9\linewidth}{\centering
    \PH{Abbildung: Micro-F$_1$ (y) gegen Parameterzahl des Textmodells (x,
    logarithmisch) für DeBERTa-v3 (Fine-Tuning), Qwen3.5 (Zero-Shot) und
    Qwen3.5 (QLoRA); eine Linie je Familie und Regime, Fehlerbalken =
    Standardabweichung über Seeds}}}
  \caption{Erkennungsleistung in Abhängigkeit von der Parameterzahl,
  getrennt nach Modellfamilie und Regime.
  Parameterzahlen nach Tabelle~\ref{tab:experiment-matrix}
  (DeBERTa-v3: 184 und 434~Mio.; Qwen3.5: 0{,}75, 4{,}21 und 26{,}90~Mrd.).}
  \label{fig:groesse}
\end{figure}

Innerhalb der Encoder-Familie bringt der Schritt von 184 auf 434~Mio.
Parameter \PH{+x,x} Punkte.
Innerhalb der LLM-Familie wächst der Micro-F$_1$ im QLoRA-Regime von
\PH{xx,x} (0.8B) über \PH{xx,x} (4B) auf \PH{xx,x} (27B); der Zugewinn von
0.8B auf 4B beträgt \PH{x,x} Punkte, der von 4B auf 27B \PH{x,x} Punkte.
\PH{Einordnen: Sättigt der Zugewinn ab 4B? Ist der Schritt 4B zu 27B
überhaupt unterscheidbar (Summe der SD)?}
Im Zero-Shot-Regime fällt der Größeneffekt \PH{deutlich größer / ähnlich}
aus: \PH{xx,x}, \PH{xx,x} und \PH{xx,x} Punkte.
\PH{Falls das 0.8B-Modell im Zero-Shot nahezu versagt: das als Befund
festhalten (Formtreue in Abschnitt~\ref{sec:erg-formtreue} dazu heranziehen).}
Über die Familien hinweg erreicht \PH{das kleinste LLM im QLoRA-Regime
(0,75 Mrd.) / DeBERTa-v3-base (184 Mio.)} \PH{xx,x} Punkte und damit
\PH{mehr / weniger} als \PH{DeBERTa-v3-base / Qwen3.5-0.8B}, obwohl es
\PH{viermal so viele / ein Viertel so viele} Parameter besitzt.
Mit zwei Encoder- und drei LLM-Größen bleibt diese Betrachtung
deskriptiv (Abschnitt~\ref{sec:einfluss-parameteranzahl}); zudem
variieren Hyperparameter und GPU-Typ mit der Größe.

\section{Effizienzanalyse}
\label{sec:erg-effizienz}

% TODO: inference_metrics.yaml je Lauf: Latenz (Mittel, Median, P95),
% Durchsatz (LLMs), max. allokierter GPU-Speicher (MiB -> GB);
% results.yaml: Trainingsdauer (27B nur Untergrenze, Fußnote).
% Mittelwerte über Seeds; Abbildung fig:tradeoff (F1 gegen Latenz, F1
% gegen Speicher) aus tab:haupt und tab:effizienz erzeugen.

Tabelle~\ref{tab:effizienz} stellt den Ressourcenbedarf der acht
Konfigurationen nach dem Protokoll aus
Abschnitt~\ref{sec:effizienzmetriken} zusammen.
Die Latenz bezeichnet die mittlere Dauer des Modellaufrufs je Testsatz
bei Batchgröße~1; der GPU-Speicher ist der maximal von PyTorch
allokierte Speicher während der Testinferenz.
Da die Konfigurationen auf drei GPU-Typen liefen
(Tabelle~\ref{tab:hardware}), ist der GPU-Typ in der Tabelle
mitgeführt; Vergleiche werden vorrangig innerhalb gleicher Hardware
gezogen.

\begin{table}[htbp]
  \centering
  \footnotesize
  \caption{Ressourcenbedarf der Konfigurationen.
  Training: Wall-Clock-Dauer des Trainingsaufrufs in Stunden, Mittel über
  drei Seeds; Latenz: Dauer des Modellaufrufs je Testsatz bei
  Batchgröße~1 in Millisekunden (Mittelwert, Median, 95.~Perzentil);
  Durchsatz: Testsätze je Sekunde reiner Generierungszeit; Speicher:
  maximal allokierter GPU-Speicher während der Testinferenz in GB.
  Parameter: Textmodell gesamt / trainierbar
  (Tabelle~\ref{tab:experiment-matrix}).
  GPU: H100 = H100 80\,GB HBM3, A100 = A100-SXM4 80\,GB, NVL = H100 NVL
  94\,GB.}
  \label{tab:effizienz}
  \setlength{\tabcolsep}{3pt}
  \begin{tabular}{llrrrrrrr}
    \toprule
    \textbf{Modell} & \textbf{Regime} & \textbf{GPU} &
    \textbf{Training} & \multicolumn{3}{c}{\textbf{Latenz [ms]}} &
    \textbf{Durchsatz} & \textbf{Speicher} \\
    \cmidrule(lr){5-7}
    & & & \textbf{[h]} & \textbf{Mittel} & \textbf{Median} & \textbf{P95} &
    \textbf{[Sätze/s]} & \textbf{[GB]} \\
    \midrule
    DeBERTa-v3-base  & Fine-Tuning & H100 & \PH{x,x} & 8{,}7 & \PH{x,x} & \PH{x,x} & -- & \PH{x,x} \\
    DeBERTa-v3-large & Fine-Tuning & H100 & \PH{x,x} & \PH{xx,x} & \PH{x,x} & \PH{x,x} & -- & \PH{x,x} \\
    \midrule
    Qwen3.5-0.8B & Zero-Shot & A100 & -- & \PH{xxx} & \PH{xxx} & \PH{xxx} & \PH{x,x} & \PH{x,x} \\
    Qwen3.5-4B   & Zero-Shot & A100 & -- & \PH{xxx} & \PH{xxx} & \PH{xxx} & \PH{x,x} & \PH{xx,x} \\
    Qwen3.5-27B  & Zero-Shot & NVL  & -- & \PH{xxxx} & \PH{xxxx} & \PH{xxxx} & \PH{x,x} & \PH{xx,x} \\
    \midrule
    Qwen3.5-0.8B & QLoRA & H100 & \PH{x,x} & \PH{xxx} & \PH{xxx} & \PH{xxx} & \PH{x,x} & \PH{x,x} \\
    Qwen3.5-4B   & QLoRA & A100 & \PH{x,x} & \PH{xxx} & \PH{xxx} & \PH{xxx} & \PH{x,x} & \PH{x,x} \\
    Qwen3.5-27B  & QLoRA & NVL  & $\geq$\,\PH{xx,x}$^{a}$ & \PH{xxxx} & \PH{xxxx} & \PH{xxxx} & \PH{x,x} & \PH{xx,x} \\
    \bottomrule
    \multicolumn{9}{@{}l}{\scriptsize $^{a}$ Nur der Abschnitt nach der
      Wiederaufnahme protokolliert; Untergrenze
      (Abschnitt~\ref{sec:effizienzmetriken}).} \\
  \end{tabular}
\end{table}

\paragraph{Trainingsdauer.}
Das vollständige Fine-Tuning der Encoder dauerte \PH{x,x}~h (base,
fünf Epochen) bzw. \PH{x,x}~h (large, vier Epochen).
Das QLoRA-Training über drei Epochen benötigte \PH{x,x}~h für das
0.8B-Modell und \PH{x,x}~h für das 4B-Modell; für das 27B-Modell ist nur
eine Untergrenze von \PH{xx}~h bekannt.
\PH{Einordnen: Verhältnis der Trainingsdauern LLM zu Encoder; Hinweis,
dass die LLM-Läufe generative Validierung einschließen und das 4B-Modell
auf einer A100 lief.}

\paragraph{Inferenzlatenz.}
Die Encoder verarbeiten einen Testsatz in \PH{8{,}7} ms (base) bzw.
\PH{xx,x} ms (large).
Die LLMs benötigen im QLoRA-Regime \PH{xxx} ms (0.8B), \PH{xxx} ms (4B) und
\PH{xxxx} ms (27B) je Satz, also das \PH{xx}-, \PH{xx}- bzw. \PH{xxx}-Fache
des schnelleren Encoders.
\PH{Falls das 0.8B-Modell und DeBERTa-v3-base auf derselben H100 liefen:
diesen Vergleich als hardwaregleich hervorheben.}
Das 95.~Perzentil liegt bei den LLMs \PH{um das x-Fache} über dem Median,
was auf die Abhängigkeit der Generierungsdauer von der Zahl der
ausgegebenen Entitäten zurückgeht (Abschnitt~\ref{sec:effizienzmetriken}).
Die Zero-Shot-Läufe sind \PH{langsamer / schneller} als die QLoRA-Läufe
derselben Größe; sie unterscheiden sich in Präzision (BF16 statt NF4),
Decoding (Sampling statt Greedy) und beim 0.8B-Modell im GPU-Typ, sodass
dieser Unterschied nicht einem Faktor zugeordnet werden kann.

\paragraph{GPU-Speicher.}
Der Spitzenspeicher der Inferenz beträgt \PH{x,x} GB (DeBERTa-v3-base),
\PH{x,x} GB (DeBERTa-v3-large), \PH{x,x} GB (Qwen3.5-0.8B, QLoRA),
\PH{x,x} GB (Qwen3.5-4B, QLoRA) und \PH{xx,x} GB (Qwen3.5-27B, QLoRA).
Im Zero-Shot-Regime mit BF16-Gewichten liegt der Bedarf bei \PH{x,x},
\PH{xx,x} und \PH{xx,x} GB; der Unterschied zum QLoRA-Regime entspricht
im Wesentlichen der 4-Bit-Quantisierung des Basismodells.
\PH{Einordnen: Welche Konfigurationen passen auf eine GPU mit 24 GB bzw.
mit 16 GB? Das ist die Grundlage für FF3.}

\paragraph{Trade-off.}
Abbildung~\ref{fig:tradeoff} stellt den Micro-F$_1$ der Latenz und dem
GPU-Speicher gegenüber.
\PH{Beschreiben: Liegt eine Konfiguration auf der Pareto-Front (kein
anderes Modell ist zugleich besser und sparsamer)? Typischerweise: die
beiden Encoder und das beste LLM; einordnen, ob das 4B- oder das
27B-Modell dominiert wird.}

\begin{figure}[htbp]
  \centering
  \fbox{\parbox[c][6cm][c]{0.9\linewidth}{\centering
    \PH{Abbildung mit zwei Teilbildern: links Micro-F$_1$ (y) gegen
    mittlere Inferenzlatenz in ms (x, logarithmisch); rechts Micro-F$_1$
    gegen GPU-Spitzenspeicher in GB (x, logarithmisch). Ein Punkt je
    Konfiguration, Form nach Regime, Farbe nach Familie, GPU-Typ als
    Beschriftung; Pareto-Front als Linie}}}
  \caption{Erkennungsleistung gegenüber Inferenzlatenz (links) und
  GPU-Spitzenspeicher (rechts).
  Latenz und Speicher nach Tabelle~\ref{tab:effizienz}; Werte verschiedener
  GPU-Typen sind nur eingeschränkt vergleichbar.}
  \label{fig:tradeoff}
\end{figure}

\section{Fehleranalyse}
\label{sec:erg-fehlertypen}

% TODO: aus test_predictions.json je Lauf: Referenz- und vorhergesagte
% Spans aus den BIO-Sequenzen extrahieren (seqeval-Logik, strikt, IOB2)
% und jeden Referenzspan einer Klasse zuordnen:
%   korrekt (Grenzen und Typ gleich), Typfehler (Grenzen gleich, Typ
%   anders), Grenzfehler (überlappend, Typ gleich), Grenz- und Typfehler
%   (überlappend, Typ anders), Auslassung (keine Überlappung).
% Vorhergesagte Spans ohne Überlappung = Zusätze. Mittel über Seeds.
% Beispiele je Fehlerklasse für tab:app-beispiele (Anhang A.2) auswählen.

Um die Fehlerprofile der Modellklassen zu vergleichen, wird jede
Referenzentität nach dem in Abschnitt~\ref{sec:performanzmetriken}
festgelegten Schema genau einer Klasse zugeordnet: korrekt erkannt,
Typfehler, Grenzfehler, Grenz- und Typfehler oder Auslassung;
vorhergesagte Spans ohne Überlappung mit einer Referenzentität zählen als
Zusätze.
Tabelle~\ref{tab:fehlertypen} zeigt die Verteilung.

\begin{table}[htbp]
  \centering
  \footnotesize
  \caption{Fehlerprofile auf dem Testsplit: Verteilung der 25\,645
  Referenzentitäten auf Fehlerklassen (in Prozent) sowie Zahl der
  Zusätze (vorhergesagte Spans ohne Überlappung mit einer Referenz).
  Trainierte Konfigurationen als Mittel über drei Seeds.}
  \label{tab:fehlertypen}
  \setlength{\tabcolsep}{3pt}
  \begin{tabular}{llrrrrrr}
    \toprule
    \textbf{Modell} & \textbf{Regime} & \textbf{Korrekt} &
    \textbf{\shortstack[c]{Typ-\\fehler}} & \textbf{\shortstack[c]{Grenz-\\fehler}} &
    \textbf{\shortstack[c]{Grenz- und\\Typfehler}} &
    \textbf{\shortstack[c]{Aus-\\lassung}} & \textbf{Zusätze} \\
    \midrule
    DeBERTa-v3-base  & Fine-Tuning & \PH{xx,x} & \PH{x,x} & \PH{x,x} & \PH{x,x} & \PH{x,x} & \PH{xxx} \\
    DeBERTa-v3-large & Fine-Tuning & \PH{xx,x} & \PH{x,x} & \PH{x,x} & \PH{x,x} & \PH{x,x} & \PH{xxx} \\
    Qwen3.5-0.8B & QLoRA & \PH{xx,x} & \PH{x,x} & \PH{x,x} & \PH{x,x} & \PH{x,x} & \PH{xxx} \\
    Qwen3.5-4B   & QLoRA & \PH{xx,x} & \PH{x,x} & \PH{x,x} & \PH{x,x} & \PH{x,x} & \PH{xxx} \\
    Qwen3.5-27B  & QLoRA & \PH{xx,x} & \PH{x,x} & \PH{x,x} & \PH{x,x} & \PH{x,x} & \PH{xxx} \\
    Qwen3.5-27B  & Zero-Shot & \PH{xx,x} & \PH{x,x} & \PH{x,x} & \PH{x,x} & \PH{x,x} & \PH{xxxx} \\
    \bottomrule
  \end{tabular}
\end{table}

\PH{Beschreiben: Welche Fehlerklasse dominiert bei den Encodern, welche
bei den LLMs? Zhan et al. (2026) fanden auf CoNLL-2003, dass der Encoder
vor allem Entitäten übersieht und das LLM vor allem falsche Typen
vergibt; prüfen, ob sich das hier wiederholt. Grenzfehler bei
mehrteiligen Namen (ORG, MEDIA) und Typfehler zwischen ANIM/PLANT/FOOD
sowie LOC/ORG gesondert nennen.}
Beispiele typischer Fehler jeder Klasse zeigt
Tabelle~\ref{tab:app-beispiele} im Anhang.

\section{Transfer auf CrossNER}
\label{sec:erg-crossner}

% TODO: results/crossner/<lauf>/<domäne>/: P/R/F1 je Domäne und Seed für
% DeBERTa-v3-large und Qwen3.5-27B (QLoRA); gepoolter Mikrowert und
% Domänenmittel je Seed -> Mittel +- SD; MultiNERD-Referenz (nur
% PER/ORG/LOC) aus den MultiNERD-Testvorhersagen derselben Checkpoints;
% Sensitivität ohne die 21 überschneidenden Sätze; Parse-Fehlerrate,
% Latenz, Speicher auf CrossNER (beide auf H100 80 GB).

Tabelle~\ref{tab:crossner-erg} berichtet den Ergänzungsversuch aus
Abschnitt~\ref{sec:crossner}: die jeweils größten Modelle beider
Familien, auf MultiNERD trainiert, ausgewertet auf den Testsplits der
fünf CrossNER-Domänen und beschränkt auf die Grobtypen PER, ORG und LOC.

\begin{table}[htbp]
  \centering
  \footnotesize
  \caption{CrossNER-Transfertest (nur PER, ORG, LOC; strikte Werte in
  Punkten; Mittelwert~$\pm$~Standardabweichung über die drei
  MultiNERD-Checkpoints je Modell).
  Gepoolt: Mikrowert über alle 2\,506 Sätze; Domänenmittel: ungewichtetes
  Mittel der fünf Domänen-F$_1$.
  MultiNERD-Referenz: dieselben Checkpoints auf dem MultiNERD-Testsplit,
  ebenfalls auf PER, ORG und LOC beschränkt; $\Delta$F$_1$ = gepoolter
  CrossNER-Wert minus Referenz.
  Sensitivität: gepoolter Wert ohne die 21 Sätze, die mit
  MultiNERD-Trainingssätzen übereinstimmen.}
  \label{tab:crossner-erg}
  \setlength{\tabcolsep}{3pt}
  \begin{tabular}{>{\raggedright\arraybackslash}p{3.2cm}rcccccc}
    \toprule
    & & \multicolumn{3}{c}{\textbf{DeBERTa-v3-large}} &
      \multicolumn{3}{c}{\textbf{Qwen3.5-27B (QLoRA)}} \\
    \cmidrule(lr){3-5}\cmidrule(lr){6-8}
    \textbf{Domäne} & \textbf{Sätze} & \textbf{P} & \textbf{R} & \textbf{F$_1$} &
      \textbf{P} & \textbf{R} & \textbf{F$_1$} \\
    \midrule
    Künstliche Intelligenz & 431 & \PH{xx,x} & \PH{xx,x} & \PH{xx,x}\,$\pm$\,\PH{x,x} & \PH{xx,x} & \PH{xx,x} & \PH{xx,x}\,$\pm$\,\PH{x,x} \\
    Literatur              & 416 & \PH{xx,x} & \PH{xx,x} & \PH{xx,x}\,$\pm$\,\PH{x,x} & \PH{xx,x} & \PH{xx,x} & \PH{xx,x}\,$\pm$\,\PH{x,x} \\
    Musik                  & 465 & \PH{xx,x} & \PH{xx,x} & \PH{xx,x}\,$\pm$\,\PH{x,x} & \PH{xx,x} & \PH{xx,x} & \PH{xx,x}\,$\pm$\,\PH{x,x} \\
    Politik                & 651 & \PH{xx,x} & \PH{xx,x} & \PH{xx,x}\,$\pm$\,\PH{x,x} & \PH{xx,x} & \PH{xx,x} & \PH{xx,x}\,$\pm$\,\PH{x,x} \\
    Naturwissenschaften    & 543 & \PH{xx,x} & \PH{xx,x} & \PH{xx,x}\,$\pm$\,\PH{x,x} & \PH{xx,x} & \PH{xx,x} & \PH{xx,x}\,$\pm$\,\PH{x,x} \\
    \midrule
    Gepoolt                & 2\,506 & \PH{xx,x} & \PH{xx,x} & \PH{xx,x}\,$\pm$\,\PH{x,x} & \PH{xx,x} & \PH{xx,x} & \PH{xx,x}\,$\pm$\,\PH{x,x} \\
    Domänenmittel          & -- & -- & -- & \PH{xx,x}\,$\pm$\,\PH{x,x} & -- & -- & \PH{xx,x}\,$\pm$\,\PH{x,x} \\
    Sensitivität (ohne 21 Sätze) & 2\,485 & -- & -- & \PH{xx,x} & -- & -- & \PH{xx,x} \\
    \midrule
    MultiNERD-Referenz (PER/ORG/LOC) & 16\,454 & \PH{xx,x} & \PH{xx,x} & \PH{xx,x} & \PH{xx,x} & \PH{xx,x} & \PH{xx,x} \\
    $\Delta$F$_1$ (gepoolt $-$ Referenz) & & & & \PH{$-$xx,x} & & & \PH{$-$xx,x} \\
    \bottomrule
  \end{tabular}
\end{table}

Gegenüber der MultiNERD-Referenz fällt der F$_1$ auf CrossNER bei
DeBERTa-v3-large um \PH{xx,x} Punkte und bei Qwen3.5-27B um \PH{xx,x}
Punkte.
\PH{Beschreiben: Welches Modell verliert weniger? Ist der Unterschied
zwischen den Modellen größer als die Streuung? Welche Domäne ist am
schwersten (vermutlich Naturwissenschaften oder KI wegen vieler
nicht abgebildeter Typen, die als Zusätze zählen) und welche am
leichtesten (vermutlich Politik mit vielen PER/ORG/LOC)? Precision
gegenüber Recall: Verlieren die Modelle vor allem Recall, weil
CrossNER-Typen wie politician auf PER abgebildet werden, oder
Precision, weil misc-Spans als PER/ORG/LOC vorhergesagt werden?}
Der Ausschluss der 21 mit MultiNERD-Trainingssätzen übereinstimmenden
Sätze verändert die gepoolten Werte um \PH{höchstens x,x} Punkte.
Beim LLM waren \PH{x,x}\,\% der Antworten nicht parsebar, \PH{nahezu alle}
davon ohne schließende Klammer; da CrossNER-Sätze \PH{länger sind und
mehr Entitäten enthalten als MultiNERD-Sätze}, wirkt das Limit von 256
Ausgabetokens hier stärker (Abschnitt~\ref{sec:crossner}), und Recall
und F$_1$ des LLM sind entsprechend unterschätzt.
Die Latenz auf CrossNER betrug \PH{xx,x} ms (DeBERTa-v3-large) und
\PH{xxxx} ms (Qwen3.5-27B) je Satz auf derselben H100.
% TODO: Satz zur 1024-Token-Sensitivitätsanalyse nur aufnehmen, wenn sie
% durchgeführt wurde; andernfalls den entsprechenden Absatz in 4.6
% streichen.

\section{Diskussion der Ergebnisse}
\label{sec:erg-diskussion}

Die folgenden Abschnitte ordnen die Befunde entlang der
Forschungsfragen aus Abschnitt~\ref{sec:problem_ziel} ein und
vergleichen sie mit den verwandten Arbeiten aus
Kapitel~\ref{chap:verwandte_arbeiten}.
Abschnitt~\ref{sec:disk-validitaet} benennt, welche Unterschiede
angesichts der Streuung und der bekannten Abweichungen
(Tabelle~\ref{tab:abweichungen}) belastbar sind.

\subsection{Erkennungsleistung im Vergleich zur Literatur}
\label{sec:disk-leistung}

% TODO: Richtung des Hauptbefunds eintragen und den passenden Absatz
% ausformulieren; den nicht zutreffenden Absatz streichen.

Der zentrale Befund zu FF1 lautet: \PH{Das beste LLM (Qwen3.5-27B, QLoRA)
liegt um x,x Punkte vor / hinter DeBERTa-v3-large; der Abstand ist
unterscheidbar / nicht unterscheidbar.}

\PH{Falls das LLM vorn liegt:} Das Ergebnis entspricht dem Muster, das
\textcite{obeidat2025llm} auf vier von fünf biomedizinischen Datensätzen
berichten, dort mit 2 bis 8 Punkten Vorsprung, und dem von
\textcite{zhan2026generativeNER} auf CoNLL-2003 mit 93{,}85 gegenüber
93{,}04 für BERT-MRC.
Der hier gemessene Vorsprung von \PH{x,x} Punkten liegt \PH{am unteren
Rand / innerhalb / oberhalb} dieser Spanne.
Anders als bei Obeidat et al., deren Vorsprung vor allem aus der
Precision stammt, geht er hier \PH{auf Precision / Recall / beides}
zurück.
\PH{Falls der Encoder vorn liegt oder kein Unterschied besteht:} Das
Ergebnis weicht von \textcite{obeidat2025llm} ab, deckt sich aber mit dem
Befund von \textcite{zhan2026generativeNER}, dass LoRA-angepasste LLMs
spezialisierte Modelle auf flachen Datensätzen nur knapp erreichen, und
mit der Beobachtung von \textcite{lu2024llmstruggle}, dass ein
LoRA-angepasstes 7B-Modell hinter einem feinjustierten Encoder bleibt.
Zwei Eigenschaften des Setups sprechen dafür, dass die Encoder hier
begünstigt sind: Die Encoder wurden vollständig feinjustiert, die LLMs
nur über Adapter mit 0{,}5 bis 0{,}8\,\% trainierbarer Parameter
(Tabelle~\ref{tab:lora-hyperparams}); und der Trainingssplit ist mit
131\,280 Sätzen groß genug, um den Vorteil des Vorwissens großer Modelle
zu verkleinern, den \textcite{bogdanov2024nuner} und
\textcite{zaratiana2024gliner} vor allem bei wenigen Beispielen finden.

Unabhängig von der Richtung ist der Abstand zwischen den trainierten
Klassen \PH{klein} im Verhältnis zum Abstand zwischen den Regimen:
Die Anpassung mit QLoRA hebt den F$_1$ um \PH{xx bis xx} Punkte über das
Zero-Shot-Regime.
\PH{Das bestätigt (prüfen: Zero-Shot deutlich unter allen trainierten
Konfigurationen) für ein feingranulares Inventar mit 15 Typen, was
\textcite{keraghel2024recent} mit GPT-4o (65{,}17 gegenüber 90{,}98 auf
CoNLL-2003) und \textcite{tan2023damo} mit ChatGPT auf MultiCoNER II
(14{,}78 gegenüber 92{,}77) für geprompte Modelle berichten: Ohne
aufgabenspezifische Anpassung bleiben LLMs weit hinter überwachten
Systemen zurück.}
Die hier gemessenen Zero-Shot-Werte von \PH{xx bis xx} Punkten liegen
\PH{deutlich unter} dem Zero-Shot-Wert von 76{,}0, den
\textcite{sainz2024gollie} für GoLLIE-7B auf einer MultiNERD-Fassung
berichten; GoLLIE ist jedoch auf Annotationsrichtlinien vieler
Datensätze trainiert und hatte acht der 15 Typen bereits gesehen, sodass
der Vergleich nur die Größenordnung einordnet.
Die trainierten Konfigurationen erreichen Werte \PH{in der Größenordnung}
der 93{,}73 (UniNER-7B) und 91{,}25 (BERT-base), die
\textcite{zhou2024universalner} auf einer MultiNERD-Variante mit
eigener Datenaufbereitung berichten; auch dort ist der Abstand zwischen
LLM und Encoder klein.
Da alle Werte dieser Arbeit die Übereinstimmung mit einer
Silberreferenz messen (Abschnitt~\ref{sec:datenauswahl}), bleibt offen,
ob die Unterschiede bei den verrauschten Typen ANIM, BIO, FOOD und PLANT
Modellfehler oder Annotationsfehler abbilden; \textcite{tedeschi2022multinerd}
berichten für diese Typen auf ihrer manuell geprüften Stichprobe
F$_1$-Werte der Annotation zwischen 0{,}52 und 0{,}71.

\subsection{Ausgabeformat und Formtreue}
\label{sec:disk-format}

Die zweite Hälfte von FF1 betrifft die Formtreue.
Im QLoRA-Regime halten die Modelle das JSON-Format mit Wortindizes
\PH{nahezu vollständig} ein: Parse-Fehlerraten von \PH{unter x}\,\%,
\PH{keine} unbekannten Typen und \PH{nur x} Text-Abweichungen auf 16\,454
Antworten.
Das stützt die in Abschnitt~\ref{sec:llm_ner_related} formulierte
Erwartung, dass Indizes, die explizit in der Eingabe stehen, den von
\textcite{zhan2026generativeNER} für Zeichenoffsets beobachteten
Einbruch vermeiden; der in Tabelle~\ref{tab:haupt} erreichte F$_1$ liegt
\PH{weit über} den rund 54 Punkten, die dort für JSON mit Zeichenoffsets
berichtet werden, und \PH{nahe an} den Inline-Formaten.
\textcite{semin2026spanLabeling} hatten die Wortindex-Variante nicht
evaluiert; die vorliegenden Ergebnisse zeigen, dass sie nach einem
Fine-Tuning \PH{tragfähig} ist.
Zugleich bestätigt sich der von \textcite{dagdelen2024structuredIE}
berichtete Befund, dass feinjustierte Modelle JSON fast immer parsebar
ausgeben, auch ohne erzwungene Formatierung.
Die verbleibenden Parse-Fehler gehen \PH{fast ausschließlich} auf das
Ausgabelimit zurück, nicht auf das Format; sie sind eine Eigenschaft der
Pipeline-Einstellung und würden mit einem höheren Limit
\PH{voraussichtlich} entfallen.

Im Zero-Shot-Regime ist die Formtreue \PH{deutlich schwächer}:
\PH{zusammenfassen, welche Fehlerklassen dominieren}.
Das zeigt, dass Formatvorgabe und Typinventar ohne Training zum
Flaschenhals werden, wie es \textcite{lu2024llmstruggle} für BIO-Tags
beschreiben; die Nachverarbeitung aus Abschnitt~\ref{sec:postprocessing}
rettet gültige Elemente, repariert aber keine fehlerhaften Indizes.
Ob Constrained Decoding diese Fehler ohne Leistungsverlust beseitigen
würde, bleibt offen; \textcite{tam2024speakfreely} berichten
Leistungseinbußen durch Formatzwang, \textcite{semin2026spanLabeling}
dagegen Gewinne durch ihr LogitMatch-Verfahren.
Für die Effizienz ist das Format ein Faktor: Weil die Ausgabelänge nur
von der Zahl der Entitäten abhängt, nicht von der Satzlänge, bleiben die
Latenzen \PH{unter / bei} denen, die \textcite{obeidat2025llm} für
BIO-Ausgaben berichten (40- bis 220-fache Encoder-Latenz); hier beträgt
der Faktor \PH{xx bis xxx}.

\subsection{Modellgröße}
\label{sec:disk-groesse}

Zu FF2 zeigen die Ergebnisse \PH{kleine / deutliche} Größeneffekte
innerhalb der Familien.
\PH{Bei den Encodern: base gegenüber large.}
Das deckt sich mit \textcite{keraghel2024recent}, bei denen
DeBERTa-v3-base und -large je nach Datensatz die Plätze tauschen.
\PH{Bei den LLMs im QLoRA-Regime: Zugewinn 0.8B zu 4B und 4B zu 27B.}
Das Muster entspricht den Beobachtungen bei GoLLIE und UniversalNER,
wo größere Backbones nur wenig hinzugewinnen und der Gewinn durch die
Anpassung den Gewinn durch die Größe übersteigt
\cite{sainz2024gollie,zhou2024universalner}.
Im Zero-Shot-Regime ist der Größeneffekt \PH{größer}, was zu der
Beobachtung passt, dass Fähigkeiten wie das Befolgen komplexer
Formatvorgaben erst ab einer gewissen Größe zuverlässig auftreten
\cite{wei2022emergent}.
Der Vergleich über die Familien hinweg \PH{zeigt, dass die Parameterzahl
allein die Leistung nicht erklärt (prüfen: gilt, wenn DeBERTa-v3-base
Qwen3.5-0.8B erreicht oder übertrifft)}: \PH{DeBERTa-v3-base mit 184 Mio.
Parametern erreicht xx,x Punkte, Qwen3.5-0.8B mit dem Vierfachen im
QLoRA-Regime xx,x.}
Architektur, Vortrainingsziel und Anpassungsverfahren überlagern den
Größenunterschied, wie \textcite{zaratiana2024gliner} und
\textcite{bogdanov2024nuner} für Encoder-Generalisten zeigen.
Die Skalierungsgesetze aus Abschnitt~\ref{sec:model_size} sind auf
diese Befunde nicht übertragbar: Sie beschreiben den Vortrainingsverlust
eigens trainierter Modellreihen, hier wurden veröffentlichte
Checkpoints mit größenabhängigen Hyperparametern und auf verschiedenen
GPUs angepasst.

\subsection{Effizienz und praktische Empfehlungen}
\label{sec:disk-effizienz}

FF3 fragt nach dem günstigsten Verhältnis von Erkennungsleistung und
Ressourcenbedarf unter Speicher- und Latenzbeschränkungen.
Die Effizienzwerte lassen sich zu drei Aussagen verdichten.
Erstens liegen die Encoder bei Latenz und Speicher um \PH{ein bis zwei}
Größenordnungen unter den LLMs: \PH{8,7 ms gegenüber xxx bis xxxx ms;
x,x GB gegenüber x,x bis xx,x GB}.
Zweitens kostet der Leistungsunterschied zwischen den Klassen
\PH{x,x} Punkte, der Effizienzunterschied dagegen den Faktor \PH{xx}
bei der Latenz; dieses Verhältnis entspricht der Schlussfolgerung von
\textcite{obeidat2025llm}, dass Encoder bei kleinen
Leistungsunterschieden oder bei Latenzanforderungen die geeignetere
Wahl bleiben.
Drittens ist das Adaptertraining selbst sparsam: Die QLoRA-Läufe
trainieren 6{,}4 bis 159~Mio. Parameter und passen mit NF4-Basismodell
auf eine GPU, das 27B-Modell benötigt dafür aber \PH{xx} GB während der
Inferenz und \PH{mindestens xx} Stunden Training.

Daraus ergeben sich Empfehlungen, die in Abschnitt~\ref{sec:fazit-zusammenfassung}
zusammengefasst werden:
\PH{(1) Bei Latenzanforderungen im Millisekundenbereich oder GPUs unter
x GB: DeBERTa-v3-base bzw. -large, mit einem Verlust von höchstens x,x
Punkten gegenüber dem besten LLM.
(2) Wenn der Leistungsunterschied von x,x Punkten den Mehraufwand
rechtfertigt und xx GB GPU-Speicher sowie Latenzen von rund x Sekunden
je Satz akzeptabel sind: Qwen3.5-27B mit QLoRA.
(3) Qwen3.5-4B mit QLoRA als Mittelweg, falls es DeBERTa-v3-large
erreicht; andernfalls dominiert.
(4) Zero-Shot nur, wenn keine Trainingsdaten vorliegen; bei 15 Typen
mit xx Punkten nicht konkurrenzfähig.}
Diese Empfehlungen gelten für die hier gemessenen Bedingungen:
Batchgröße~1, ein Satz je Anfrage, die genannten GPU-Typen.
Mit Batching steigt der Durchsatz der LLMs stärker als der der Encoder,
sodass sich die Faktoren bei Offline-Verarbeitung verschieben; dieser
Fall wurde nicht gemessen (Abschnitt~\ref{sec:abgrenzung}).

\subsection{Belastbarkeit der Befunde}
\label{sec:disk-validitaet}

Welche der berichteten Unterschiede tragen, bestimmt sich aus der
Streuung und aus den Abweichungen in Tabelle~\ref{tab:abweichungen}.
Die Standardabweichungen über die Seeds liegen bei \PH{x,x bis x,x}
Punkten; Unterschiede unterhalb von \PH{etwa x} Punkten werden daher
nicht interpretiert.
Davon betroffen sind \PH{der Vergleich DeBERTa-v3-base gegen -large /
der Schritt von 4B auf 27B / der Abstand zwischen DeBERTa-v3-large und
Qwen3.5-27B}.
Belastbar sind dagegen \PH{der Abstand zwischen Zero-Shot- und
QLoRA-Regime, der Abstand zwischen den trainierten Modellen und dem
kleinsten LLM sowie die Effizienzfaktoren}, die um ein Vielfaches über
der Streuung liegen.
Für die Effizienzwerte gilt das nur innerhalb gleicher Hardware; der
hardwaregleiche Vergleich beschränkt sich auf \PH{DeBERTa-v3 und
Qwen3.5-0.8B (QLoRA) auf der H100}.

Die Einschränkungen des Versuchsdesigns, darunter die Silberreferenz,
der einzelne Zero-Shot-Lauf je Größe, die drei GPU-Typen und die
asymmetrische Checkpoint-Auswahl, sind in Tabelle~\ref{tab:abweichungen}
zusammengefasst und werden in Abschnitt~\ref{sec:limitationen} diskutiert.
Keine von ihnen kehrt die Richtung der belastbaren Befunde um; sie
begrenzen aber, wie genau die kleinen Unterschiede zwischen den
trainierten Konfigurationen zu beziffern sind.
